%!TEX root = ../../main.tex
\section{결론}
\label{sec:disc-conclusion}

본 논문은 킬 기록을 기준으로 구성한 교전에 대해, 사전 관측 시점의 공개 상태로 교전 전후 추정 경기 승리확률의 변화 방향을 예측하는 문제를 정의하고, 자료에서 근거를 갖는 교전 정의, 경기 승패로 학습해 동결한 평가기, 승률과 시간의 유연한 기준선에 대한 짝 대비로 답하였다. 한타와 소규모 교전 각각의 안에서 교전 전 상태의 정규화 로지스틱 모델은 그 기준선보다 경기 가중 Brier 점수를 홀드아웃 패치에서 \dBrierT{}와 \dBrierS{}만큼 낮추었고, 경기 군집 부트스트랩 구간은 0을 제외하였다. 이 이득은 기저율 상수예측의 Brier에 비해 작으며, 그 크기가 곧 발견이다. 교전 전 공개 상태는 현재 승률과 경기 시각을 빼고 나면, 교전이 경기를 어디로 움직일지에 관해 측정 가능하고 작은 양의 정보를 담고 있다.

본 논문의 중심은 그 효용이 무엇에 의존하는지를 구체적으로 제시한 데 있다. 사건 특징은 프레임 특징 위에 추가 효용을 주었고 프레임 특징의 추가 효용은 코호트별로 달랐다. 라벨은 인접한 평가기와 종점 사이에서 대체로 유지되었으나 작은 변화량에서는 예측 비교의 부호가 바뀌었고, 표적을 교환가치 라벨로 바꾸면 이득의 부호도 바뀌었다. 새 프레임이 들어온 사례에서 이득은 전체와 같은 부호였고, 평가한 외부 표본에서 한타의 우세는 관측되지 않았으며 소규모 교전의 우세는 가족 보정 수준에서 불확실하였다. 공개 상태가 제공하는 추가 예측 효용을 확인하고, 그 효용과 결과 측정이 어떤 정보 구성·평가기·종점·외부 환경에 의존하는지를 제시한 것이 이 학위논문의 기여이다.
